\documentclass[conference]{IEEEtran}
\usepackage{amsmath,amssymb}
\usepackage{graphicx}
\usepackage{booktabs}
\usepackage{url}
\usepackage[hidelinks]{hyperref}
\graphicspath{{figures/}}

\begin{document}

\title{Do Robust Counterfactuals Survive Model Multiplicity in 30-Day Readmission Prediction? DiCE vs.\ Consensus and Rashomon-Robust Counterfactuals Across Heterogeneous Model Families}

\author{\IEEEauthorblockN{Srikar [SURNAME]}
\IEEEauthorblockA{[Department], [Institution]\\
[City], [Country]}}

\maketitle

\begin{abstract}
%%ABSTRACT%%
\end{abstract}

\begin{IEEEkeywords}
counterfactual explanations, algorithmic recourse, model multiplicity, Rashomon set, hospital readmission, explainable AI
\end{IEEEkeywords}

\section{Introduction}
Counterfactual explanations (CFs) tell a user how an input would have to change for a model to give a different
prediction \cite{wachter2018}. In hospital readmission prediction, a CF for a patient flagged as high-risk reads like a
recommendation: \emph{had these care-process variables been different, the model would not have flagged this
patient}. Such recommendations are only meaningful if they describe the prediction task rather than one particular
trained model. In practice, many models reach almost the same accuracy but make different individual predictions
(\emph{predictive multiplicity} \cite{marx2020multiplicity}), and a CF that flips one of them may not flip another
\cite{pawelczyk2020multiplicity,black2022consistent}.

A recent thesis \cite{emekci2026} measured this on the Diabetes 130-US readmission data \cite{strack2014}: five XGBoost
models differing only in their random seed, DiCE \cite{mothilal2020dice} counterfactuals restricted to four
care-process features, and four stability metrics. It reported a cross-model validity of 0.65 and a feature-set
consistency of 0.51, and named other model families, robustness-aware generation and consensus counterfactuals as
future work. Robust CF methods that target validity under model change already exist
\cite{upadhyay2021roar,dutta2022robx,hamman2023,jiang2024argumentative,turbal2025ellice}, and recent benchmarks show
that their robustness depends on the kind of model change \cite{kostrzewa2026which,tan2026densityflow}. However,
these benchmarks report validity and cost only: none asks whether equally good models \emph{agree on what to
recommend}, none constrains CFs to actionable features of a clinical task, and none tests transfer across
heterogeneous model families.

This paper fills that gap with a measurement study (no new CF algorithm is proposed). We ask:
\textbf{RQ1} Does the thesis's finding hold when equally accurate models come from different families (XGBoost,
logistic regression, MLP)? \textbf{RQ2} Do consensus and Rashomon-robust (ElliCE) CFs achieve higher cross-model
validity and higher agreement on the recommended features? \textbf{RQ3} What do they cost in proximity, sparsity and
plausibility? \textbf{RQ4} Which actionable features are recommended consistently?
Our contributions are: (i)~a leak-free replication of the thesis setting; (ii)~evidence that, under multiplicity,
whether a high-risk patient has any actionable recourse at all depends on which equally accurate model is deployed;
(iii)~a like-for-like comparison of DiCE, a consensus filter and ElliCE on the four thesis metrics plus quality
metrics, both in-sample and on a held-out model family; (iv)~open, CPU-only code where every reported number is
generated from versioned result files.

\section{Related Work}
\textbf{Counterfactuals and recourse.} CFs \cite{wachter2018} and diverse CFs (DiCE) \cite{mothilal2020dice} are
widely used for tabular models; recourse work adds actionability and asks whether a feasible change exists at all
\cite{ustun2019recourse}.
\textbf{Multiplicity and robust CFs.} Predictive multiplicity \cite{marx2020multiplicity} makes CFs fragile
\cite{pawelczyk2020multiplicity,black2022consistent}. Robust generators target validity under parameter shifts
(ROAR \cite{upadhyay2021roar}), for tree ensembles (RobX \cite{dutta2022robx}), with probabilistic guarantees
\cite{hamman2023}, or over a Hessian-ellipsoid approximation of the Rashomon set of linear models and MLP last layers
(ElliCE \cite{turbal2025ellice}); see \cite{jiang2024survey} for a survey.
\textbf{Agreement across models.} Ensemble-consistent CFs \cite{artelt2022ensemble}, entropic-risk CFs with a tunable
share of models that must accept the CF \cite{noorani2025entropic}, probabilistic consensus \cite{kostrzewa2026croce}
and argumentative ensembling \cite{jiang2024argumentative,jiang2025argumentative} all seek a CF accepted by many models;
our consensus method is a simple baseline from this family, not a contribution.
\textbf{Benchmarks.} \cite{kostrzewa2026which} evaluates six robust methods under eight model-change types (MLPs only);
DensityFlow \cite{tan2026densityflow} measures validity across seven model families. Neither measures structural
agreement (which features, direction, magnitude) or applies clinical actionability constraints.
\textbf{Readmission.} Readmission models are typically ensembles \cite{yu2020readmission}; on Diabetes 130-US we found
no CF study other than \cite{emekci2026}.

\section{Methods}
\subsection{Data and models}
We use the UCI Diabetes 130-US data \cite{strack2014} (101,766 encounters), keep the first encounter per patient, drop
encounters ending in death or hospice (discharge codes 11, 13, 14, 19--21), drop \texttt{weight},
\texttt{payer\_code} and \texttt{medical\_specialty} (40--97\% missing), group ICD-9 codes into nine groups, and predict
readmission within 30 days. The data are split 80/20 at the \emph{patient} level, stratified by outcome. Categorical
features are one-hot encoded and numeric features standardised (132 inputs).
We train XGBoost \cite{chen2016xgboost}, L2-regularised logistic regression (LR) and a one-hidden-layer MLP (PyTorch, CPU),
five seeds each. Seed $s$ trains on its own bootstrap resample of the training set, because seed alone does not change a
convex LR. Hyperparameters were chosen once on a validation split of the training data, all families use class
weights, and each model's threshold is set so that it flags the same 35\% of patients as high-risk (the thesis's
operating point), so a ``flip'' means the same on every model. Models within 0.02 AUC of the best form the
multiplicity (Rashomon-style) set; 0.01 is a sensitivity analysis.

\subsection{Actionable features and cohorts}
As in the thesis, only four care-process variables may change: time in hospital, number of lab procedures, number of
procedures and number of medications, each within its training 1st--99th percentile range; every other feature is
immutable (verified by a test over all generated CFs). An exhaustive grid over these four features tells us, for every
patient and model, whether \emph{any} counterfactual exists. The \textbf{coverage cohort} is 300 random test patients
flagged by all band models; it measures how often recourse exists and how often each method finds it. The
\textbf{stability cohort} contains every test patient flagged by all band models for whom recourse exists under every
model, mirroring the thesis, which kept only patients for whom every model produced CFs.

\subsection{Counterfactual methods}
\textbf{DiCE} (random method) generates CFs on each source model. \textbf{Consensus} filters each source model's ten
DiCE candidates, keeping those that at least $\tau$ of a committee of models also classify as low-risk ($\tau=0.8$;
0.6 and 1.0 in the sensitivity analysis), and returns the closest. \textbf{ElliCE} \cite{turbal2025ellice} returns one
CF per LR or MLP source that is valid for every model in an ellipsoidal approximation of the source's $\varepsilon$-Rashomon
set. For the method comparison we keep the single closest valid CF per (patient, source) for every method, so that
DiCE's deliberate diversity is not counted as instability; DiCE and consensus select from the same candidate pool, so they
differ only in the agreement filter.

\subsection{Metrics and protocol}
We use the thesis's four metrics unchanged \cite{emekci2026}: \emph{cross-model validity} (share of evaluation models
that classify the CF as low-risk), \emph{feature-set consistency} (Jaccard overlap of the changed features over CF pairs
from different source models), \emph{direction consistency} (share of changes to a feature, made by CFs from at least two
source models, that follow the dominant sign) and \emph{magnitude consistency} ($\max(0, 1-|\Delta_a-\Delta_b|)$ for
same-direction pairs, with changes normalised by the feature's range). Quality metrics are proximity (range-normalised
L1 change), sparsity (number of changed features), plausibility (distance to the nearest training patient that the source
model predicts as low-risk) and coverage. Consensus and ElliCE use the models they are scored on, so we report two
protocols: \emph{in-sample} (all band models) and \emph{leave-one-family-out} (LOFO): CFs are built from two families
and validity is measured on the held-out family. All estimates are patient-level means with 95\% percentile bootstrap
CIs (5,000 resamples); methods are compared with paired bootstrap differences on the same patients.

%%RESULTS%%

\bibliographystyle{IEEEtran}
\bibliography{references}

\end{document}
